\subsection{The Agentic Paradigm in Industrial Intelligence}

The concept of an autonomous agent --- a computational entity that perceives its environment, maintains an internal world model, reasons about goals and constraints, and executes purposeful actions to achieve desired outcomes --- has been a central principle in artificial intelligence research. The emergence of large language model (LLM)-based agents has fundamentally altered the applicability of this paradigm to high-stakes industrial environments, making the deployment of sophisticated autonomous agents in predictive maintenance not only technically feasible but increasingly essential for sustaining competitiveness in modern manufacturing.

The agentic layer of QASAMAP receives analytical outputs from the deep learning and quantum computing layers and transforms them into operational decisions, maintenance work orders, model update triggers, and human-interpretable reports. The empirical motivation for the agentic architecture is grounded in the dataset evidence: the 69-day monitoring period reveals a dynamic operational environment in which daily anomaly rates fluctuate, machine health trajectories diverge, and maintenance demand patterns evolve continuously. A static analytical system cannot adapt to such non-stationarity without expensive manual retraining cycles. The agentic layer is the mechanism through which QASAMAP transcends static system limitations, implementing autonomous threshold calibration and continuous learning from operational feedback.

\subsection{Multi-Agent Orchestration Architecture}

The QASAMAP agentic layer implements a sequential multi-agent pipeline comprising four specialised LLM-powered agents coordinated by the \texttt{MultiAgentOrchestrator} class in \texttt{quantum\_agent.py}. All agents use the \texttt{gpt-oss:120b} model via the Ollama Cloud API. The four agents --- Monitoring, Diagnosis, Planning, and Reporting --- execute in a fixed sequence for each machine event, with each agent's structured output feeding into the next agent's prompt context. Figure~\ref{fig:agentic_ai_architecture} illustrates the complete orchestration pipeline including the \texttt{ContextMemoryStore}, \texttt{ReasoningEngine}, quantum override logic, Action Layer, and \texttt{LearningLoop}.

The pipeline is initiated per machine event after the quantum layer has run. For each machine, the orchestrator assembles a unified context record containing: the raw sensor readings; the classical risk score and risk level (LOW/MEDIUM/HIGH/CRITICAL) from the \texttt{ReasoningEngine}; the per-sensor Pauli-Z expectations $\langle Z_i \rangle$ and composite label from the \texttt{PauliZRiskEncoder}; the QUBO priority rank; and the QAOA slot assignment. This unified context is passed through all four agents in sequence.

\subsection{Agent Roles and Responsibilities}

\paragraph{Monitoring Agent.}
The Monitoring Agent is the first agent in the pipeline. It receives the sensor readings, classical risk score, and quantum context and returns a structured JSON assessment containing: \texttt{risk\_level} (LOW/MEDIUM/HIGH/CRITICAL), \texttt{anomaly\_type} (a descriptive label such as ``vibration anomaly'' or ``thermal event''), \texttt{quantum\_confirmation} (whether quantum signals corroborate the classical assessment), and \texttt{confidence} (a float in $[0,1]$). The Monitoring Agent's role is analogous to a first-line sensor analyst: it synthesises the raw evidence into an initial diagnostic hypothesis without yet committing to a maintenance action. In the quantum-augmented pipeline, the agent explicitly acknowledges the Pauli-Z quantum label and QUBO rank as secondary evidence sources, enabling it to flag concern for machines that classical thresholds alone would route to passive monitoring.

\paragraph{Diagnosis Agent.}
The Diagnosis Agent receives the Monitoring Agent's assessment together with the per-sensor Pauli-Z values as quantum evidence. It returns a structured diagnosis containing: \texttt{probable\_cause} (e.g., ``bearing wear'', ``overheating'', ``pressure drop''), \texttt{affected\_component} (e.g., ``spindle bearing'', ``cooling system''), \texttt{urgency\_hours} (an integer estimate of hours until intervention is required), and \texttt{quantum\_evidence\_used} (whether and which Pauli-Z values influenced the diagnosis). The Diagnosis Agent functions as a domain expert reasoning system: it uses the individual $\langle Z_i \rangle$ values to identify which sensor is in the quantum danger zone and constructs a mechanistic explanation linking the sensor pattern to a probable failure mode. This interpretable, sensor-level diagnosis satisfies the explainability requirements of industrial AI governance and provides maintenance technicians with actionable evidence rather than opaque model scores.

\paragraph{Planning Agent.}
The Planning Agent receives the diagnostic output together with the QUBO rank and QAOA slot assignment and determines a concrete maintenance action from four options: \texttt{monitor} (passive continuous surveillance), \texttt{inspect} (targeted physical inspection at the next opportunity), \texttt{schedule\_maintenance} (formal work order for a planned maintenance window), or \texttt{immediate\_shutdown} (emergency halt). The agent also assigns an action priority P1 (critical, $\leq$4 hours) through P4 (routine, next scheduled window) and sets a \texttt{quantum\_driven} flag indicating whether the quantum override mechanism was responsible for escalating the decision above the classical risk level. High-risk actions (\texttt{schedule\_maintenance}, \texttt{immediate\_shutdown}) are routed through the human-in-the-loop approval gate before execution.

\paragraph{Reporting Agent.}
The Reporting Agent synthesises the outputs of all three preceding agents into a structured natural language maintenance alert. The alert summarises machine identity, current sensor readings, classical and quantum risk assessments, the Diagnosis Agent's finding, the Planning Agent's recommended action and priority, and the estimated urgency window. The report is designed for direct communication to maintenance supervisors and plant managers, bridging the gap between the system's computational outputs and the human operational context in which maintenance decisions are executed.

\subsection{Context Memory and Reasoning Engine}

The \texttt{ContextMemoryStore} maintains four persistent data structures across the full pipeline execution. A \texttt{sliding\_window} implemented as a \texttt{deque} of maximum length 10 stores the most recent sensor events per machine, providing short-term temporal context for agent reasoning. A \texttt{decision\_log} records all agent actions with timestamps, risk scores, and quantum override flags, enabling retrospective audit and accountability tracking. A \texttt{feedback\_log} records post-action outcomes for the \texttt{LearningLoop} and a \texttt{quantum\_log} stores all Pauli-Z encodings, QUBO results, and QAOA assignments for traceability.

The \texttt{ReasoningEngine} computes classical risk scores via a multi-factor scoring function. Warning-level threshold crossings for temperature (82\textdegree C), vibration (55 units), energy (3.8 kWh), humidity (68\%), and pressure (4.3 bar) contribute 0.15 to the risk score; critical crossings contribute 0.35. RUL below 50 hours contributes 0.30, anomaly flag contributes 0.20, and downtime risk above 0.6 contributes 0.15. The composite score is clipped to $[0,1]$ and classified as LOW ($<0.30$), MEDIUM (0.30--0.60), HIGH (0.60--0.75), or CRITICAL ($\geq0.75$).

\subsection{Quantum Override Logic}

The \texttt{should\_quantum\_override} method in \texttt{QuantumOrchestrator} provides the key mechanism by which the quantum layer enhances the classical agentic pipeline. An override is triggered when any of three conditions holds: the machine's Pauli-Z label is \texttt{QUANTUM-ANOMALY} or \texttt{QUANTUM-WARN}; the machine was selected by the QUBO annealer as highest priority; or the machine was assigned to QAOA Slot~A (IMMEDIATE/URGENT) with quantum risk above 0.40. When triggered, the orchestrator bypasses the \texttt{auto-LOW-skip} rule --- which would otherwise route LOW-risk machines to passive monitoring without agent invocation --- and escalates the machine through the full four-agent reasoning chain. The override reason is documented in the decision log for transparency.

In the non-quantum agentic pipeline, the \texttt{auto-LOW-skip} rule causes 6 out of 7 machines in the experimental evaluation to receive passive monitoring without active agent reasoning. The quantum override mechanism raises this to active agent processing for 6 out of 7 machines, demonstrating that the quantum pre-screening layer fundamentally changes the coverage profile of the agentic system.

\subsection{Action Layer and Human-in-the-Loop Governance}

The \texttt{execute\_action} function dispatches recommended actions to simulated external operational systems: a CMMS for work order creation, an ERP system for spare parts availability checking, a supervisor notification system, and a maintenance scheduling calendar. All actions are timestamped and logged with full provenance including quantum override context and the natural language report from the Reporting Agent.

High-risk actions (\texttt{schedule\_maintenance}, \texttt{immediate\_shutdown}) require explicit supervisor approval before execution, implementing a human-in-the-loop safety gate that retains human authority over consequential decisions while delegating routine monitoring and inspection decisions to the autonomous pipeline. Override decisions by human supervisors are logged and fed back to the \texttt{LearningLoop} as labelled examples of expert judgment, enabling progressive calibration of agent thresholds to align with site-specific operational norms.

\subsection{Continuous Learning Loop}

The \texttt{LearningLoop} class implements two adaptation mechanisms that address the non-stationarity of real industrial environments. The classical feedback mechanism evaluates false positive and false negative rates from the \texttt{feedback\_log} after each decision cycle and adjusts \texttt{ReasoningEngine} sensor thresholds: if the false positive rate exceeds 40\%, temperature and vibration warning thresholds are relaxed by 2\textdegree C and 2 units respectively; if the false negative rate exceeds 20\%, the temperature warning threshold is tightened by 2\textdegree C.

The quantum feedback mechanism cross-checks Pauli-Z labels against final agent actions: machines labelled \texttt{QUANTUM-ANOMALY} or \texttt{QUANTUM-WARN} that received only passive monitoring without quantum override are flagged as potential quantum false negatives (Q-FN) and logged for human review. When the \texttt{feedback\_log} accumulates at least 5 records, a structured retraining payload is exported to \texttt{retrain\_payload.json} containing the most recent sensor readings, updated threshold parameters, and flagged Q-FN cases for submission to the T-GCN retraining pipeline, closing the closed-loop adaptation cycle.